% Standalone appendix fragment; the root manuscript supplies the section heading.
% Literature keys: chierchia2021measure, biasco2018measure.
% All angular variables in this appendix have period 2*pi.
\label{app:measure-proof}
\subsection{A fixed-domain bound and its scope}
\label{sec:measure-fixed-tube}

This appendix quantifies a sufficient local measure bound for the
Hamiltonian identified from COMSOL integral observations. The construction
uses the exact finite-amplitude torus and bounds the complete transformed
Hamiltonian, including the action-domain remainder. The resulting width
illustrates the conservatism of the explicit constants in
\citet{chierchia2021measure}.

\begin{samepage}
\begin{theorem}[Explicit finite-amplitude tubular measure]
\label{thm:measure-finite-tube}
Let $h$ be the fitted Hamiltonian certified in
Section~\ref{sec:validated-integral-learning}, and let $K_\infty$ be its
exact invariant torus. There is an explicitly constructed analytic
symplectic chart $\Xi=\Psi\circ\Phi_{\chi}$, specified in
Sections~\ref{sec:measure-explicit-tube}--\ref{sec:measure-tube-normal-form},
that fixes this torus at zero action.
For the fixed domain
\begin{equation}
 U_w=\Xi\bigl(\mathbb T^2_{2\pi}\times(-w,w)^2\bigr),
 \qquad w=10^{-220},
 \label{eq:measure-fixed-tube}
\end{equation}
the KAM fraction satisfies
\begin{equation}
 f_{\rm KAM}(U_w)\geq0.9669383.
 \label{eq:measure-certified-fraction}
\end{equation}
The centre torus is the previously certified motion with the
$1.8977904$ mm tip witness and the $1.8977962$ mm all-phase centreline
bound. The width $w$ is a dimensionless action half-width, not a
displacement amplitude.
\end{theorem}
\end{samepage}

\begin{proof}
The explicit chart and a convergent quadratic normalization give
\[
 h\circ\Xi=e+\omega\cdot I+\tfrac12 I^\top\overline A I+R_3,
 \qquad \omega=2\pi\varpi.
\]
The outward bounds are
$\|\overline A\|_\infty\leq3$,
$\|\overline A^{-1}\|_\infty\leq5100$ and
$\|R_3\|_{2w,1/128}\leq4\times10^{51}w^3$.
Use the global KAM theorem with real action radius $w$, complex radius
$2w$, final angle strip $1/256$, frequency exponent $\tau=2$,
boundary allowance $\rho=w/500$ and
$\alpha=w/(1000\times5100)$.
The dimensionless smallness quantity is at most
$3.1212\times10^{-155}$, below the theorem threshold
$7.6577\times10^{-155}$.
The geometric estimate in Section~\ref{sec:measure-ck-bound} gives
\begin{equation}
 \begin{aligned}
 1-f_{\rm KAM}(U_w)
 &\leq\frac{25}{16}\left(
      \frac{4\rho}{w}+\frac{8\alpha L\zeta(2)}{w}\right)\\
 &=\frac{25}{16}\left(0.008+0.008\zeta(2)\right)
 <0.0330617,
 \end{aligned}
 \label{eq:measure-budget}
\end{equation}
where $L=5100$. The two terms account for the boundary layer and all
integer resonance strips, respectively. Symplecticity transfers this
relative measure to $U_w$ without a density correction.
\end{proof}

\begin{table}[t]
\centering\small
\caption{Outward bounds for the explicit finite-amplitude tube.
The smallness quotient is the tested perturbation divided by the
sufficient threshold; failure of that test does not imply chaos.}
\label{tab:measure-tube-certificate}
\begin{tabular}{lr}
\hline
Quantity & Bound or value\\
\hline
Initial chart action radius & $10^{-3}$\\
Quadratic-normalization radius & $10^{-17}$\\
Certified tube half-width $w$ & $10^{-220}$\\
$\|\overline A\|_\infty$, $\|\overline A^{-1}\|_\infty$ & $3$, $5100$\\
Smallness quotient at $w=10^{-220}$ & $<0.407586$\\
Same quotient at $w=10^{-219}$ & $>4.0758$\\
Certified fraction in $U_w$ & $>96.6938\%$\\
\hline
\end{tabular}
\end{table}

The width in Table~\ref{tab:measure-tube-certificate} quantifies the
conservatism of this general theorem. The domain is centred on the exact
millimetre-amplitude torus, but its thickness is far below any operational
resolution. In particular, the computed Fourier approximation has a
certified canonical-embedding error upper bound below
$2.1\times10^{-41}$. This available guarantee cannot resolve membership
in a tube whose transverse distance is bounded by
$41w=4.1\times10^{-219}$; it is not an estimate of the actual embedding error.
The domain is mathematically specified
through $K_\infty$; this certificate does not permit an experimental
initial state to be assigned to that tube from the finite-precision
approximation. It supplies a concrete fixed-domain measure result while
leaving the much wider operating-domain question to a separate
calculation. Section~\ref{sec:measure-numerical} reports that wider-domain
numerical diagnostic, with its own error checks and interpretation.

\subsection{An explicit symplectic tube at the certified physical amplitude}
\label{sec:measure-explicit-tube}

Let $K_\infty$ be the exact torus supplied by the integral-COMSOL
certificate, with its phase normalization fixed by that construction.
Write $t=2\pi\theta$, $k(t)=K_\infty(t/(2\pi))$, and
$\omega=2\pi\varpi$. Thus $D_\omega k=X_{h_\eta}(k)$, and the
already verified Diophantine inequality becomes
\begin{equation}
 |\omega\cdot\ell|\geq\gamma_0|\ell|_1^{-3},\qquad
 \gamma_0=2\pi\frac9{1750},\quad
 0\ne\ell\in\mathbb Z^2.
 \label{eq:measure-centre-frequency}
\end{equation}
The exponent $3$ here concerns this fixed central torus. The family
measure estimate below uses the separate exponent $\tau=2$.

For each coordinate plane, define
\begin{equation}
 \begin{aligned}
 j_a(t)&=\tfrac12(k_{X_a}^2+k_{Y_a}^2),\\
 F_a(t)&=t_a+\frac{\log(1+u_{a,+})-\log(1+u_{a,-})}{2i},\\
 u_{a,\pm}(t)&=10(k_{X_a}\mp i k_{Y_a})e^{\mp it_a}-1.
 \end{aligned}
 \label{eq:measure-polar-chart-data}
\end{equation}
The logarithms use the branches at $1$. On the real torus, $F_a$ is
the phase of $k_{X_a}-ik_{Y_a}$ and $j_a>0$.
Set $B(t)=DF(t)^{-\mathsf T}$ and define
\begin{equation}
 \begin{aligned}
 J(t,I)&=j(t)+B(t)I,\\
 \Psi_{X_a}(t,I)&=\sqrt{2J_a(t,I)}\cos F_a(t),\\
 \Psi_{Y_a}(t,I)&=-\sqrt{2J_a(t,I)}\sin F_a(t).
 \end{aligned}
 \label{eq:measure-explicit-psi}
\end{equation}
Equivalently, each pair $(k_{X_a},k_{Y_a})$ is multiplied by
$\sqrt{1+(BI)_a/j_a}$; this expression is convenient for interval
evaluation and fixes the square-root branches.

This is an exact symplectic chart. Indeed, polar coordinates have
canonical one-form $J\cdot dF$, and
\[
 \Psi^*(J\cdot dF)=j(t)\cdot DF(t)dt+I\cdot dt.
\]
The first term is closed because $k$ is Lagrangian. Taking the exterior
derivative gives $\sum_a dI_a\wedge dt_a$. This argument requires
closedness, not exactness, of the first term. Since $DF$ is uniformly
close to the identity, its real lift is a global diffeomorphism of
$\mathbb T^2$; positive $J$ then makes $\Psi$ injective on its real
tube. Moreover $\Psi(t,0)=k(t)$.

The invariant-torus estimates and outward-rounded interval arithmetic
give the following upper bounds:
\begin{equation}
 \begin{gathered}
 \|k-k_0\|_{1/16}<7\,10^{-6},\\
 k_0=0.1(\cos t_1,\cos t_2,-\sin t_1,-\sin t_2),\\
 \|u_\pm\|_{1/16}<1/5000,\qquad
 \|j-0.005\mathbf1\|_{1/16}<3\,10^{-6},\\
 \|F-t\|_{1/16}<0.000201,\\
 \max\{\|DF-\mathrm{Id}\|_\infty,
             \|DF-\mathrm{Id}\|_1\}_{1/32}<0.013,\\
 \|B\|_{1/32,\infty}<1.014.
 \end{gathered}
 \label{eq:measure-chart-bounds}
\end{equation}
Here a subscript is the complex angular strip width in radians;
vector norms are maximum norms and matrix norms are induced norms.
The first bound includes the a posteriori correction to the numerical
embedding. For example, Cauchy's estimate bounds each entry of
$D(F-t)$ by $0.000201/(1/32)$, so both row and column sums are below
$0.013$. For $R_0=10^{-3}$, $|I|_\infty<R_0$ therefore implies
\[
 \begin{gathered}
 |(BI)_a/j_a|<\frac{1.014\,10^{-3}}{0.005-3\,10^{-6}}<\tfrac14,
 \\ \|\Psi\|_{1/32,R_0}<0.12.
 \end{gathered}
\]
All logarithms, square roots and inverses above are thus holomorphic
on the declared domain.

Write $h_\eta=h_2+h_{\mathrm{nl}}$. The full cosine and the exact
exported quartic coefficients give
\begin{equation}
 \|h_{\mathrm{nl}}\circ\Psi\|_{1/32,R_0}<10^{-6},\qquad
 \|h_\eta\circ\Psi\|_{1/32,R_0}<0.05.
 \label{eq:measure-chart-hamiltonian-bounds}
\end{equation}
The second estimate uses the identity $h_2\circ\Psi=\nu\cdot(j+BI)$,
whose norm is less than $7.153(0.005003+0.001014)<0.044$.
Invariance and symplecticity give the exact expansion
\begin{equation}
 h_\eta\circ\Psi=e+\omega\cdot I
                 +\tfrac12 I^\mathsf T A(t)I+O(|I|^3).
 \label{eq:measure-tube-hamiltonian}
\end{equation}
Only $h_{\mathrm{nl}}\circ\Psi$ contributes to $A$: Cauchy bounds
on its diagonal and off-diagonal second derivatives give
$\|A\|_{1/32,\infty}\leq3$.

The averaged matrix $\overline A$ is nonsingular. To check the
normalization, the tangent frame with radian angles is the certified
cycle-angle frame divided by $2\pi$; a symplectically dual normal
frame is multiplied by $2\pi$. Replacing that normal frame by
$\partial_I\Psi(t,0)$ adds a symmetric tangential shear. Its contribution
to the torsion is a directional derivative with zero angular average.
Consequently
\begin{equation}
 \overline A=(2\pi)^2\langle T_\infty\rangle,\qquad
 \|\overline A^{-1}\|_\infty
 \leq\frac{2\,10^5}{(2\pi)^2}<5100.
 \label{eq:measure-twist-normalization}
\end{equation}
The inverse-torsion bound applies to the exact final torus through
the uniform a posteriori iteration.

\subsection{A complete analytic remainder on a fixed tube}
\label{sec:measure-tube-normal-form}

Let $S$ be the symmetric, zero-average solution of
$D_\omega S=A-\overline A$. Its Fourier coefficients for $\ell\ne0$
are $S_\ell=A_\ell/(i\omega\cdot\ell)$. Equation
\eqref{eq:measure-centre-frequency} and the angular strip loss
$\delta=1/64$ imply
\begin{equation}
 \begin{aligned}
 \|S\|_{1/64,\infty}
 &\leq\frac{3}{\gamma_0}
       4\sum_{m\geq1}m^4e^{-m/64}\\
 &\leq\frac{288e^{1/64}64^5}{\gamma_0}<10^{13}=:T.
 \end{aligned}
 \label{eq:measure-homological-bound}
\end{equation}
There are $4m$ integer vectors of one-norm $m$. The second inequality
uses $\sum_{m\geq1}m^4e^{-\delta m}
\leq e^\delta\Gamma(5)\delta^{-5}$, obtained by integrating on
each interval $[m,m+1]$.

Use the exact time-one flow $\Phi_\chi$ of
$\chi(t,I)=\tfrac12 I^\mathsf TS(t)I$, not a truncated Lie series.
Fix $R_1=10^{-17}$. On the intermediate strip $3/256$, Cauchy's
estimate from the strip $1/64$ gives
$\|\partial_{t_a}S\|_\infty\leq256T$. As long as
$|I|_\infty<2R_1$, the flow satisfies
\[
 \begin{aligned}
 |\dot t|_\infty&\leq2TR_1=2\,10^{-4}<1/256,\\
 |\dot I|_\infty&\leq1024TR_1^2
       =0.1024R_1<R_1.
 \end{aligned}
\]
These strict bounds close the continuation argument from the initial
strip $1/128$ and action polydisc $R_1$: throughout the unit time
interval, the angle stays inside $3/256$ and the action inside
$2R_1<R_0$. The flow is holomorphic on the initial domain and
symplectic on its real part.

With the pullback convention $L_\chi H=\{H,\chi\}$,
$\{\omega\cdot I,\chi\}=-D_\omega\chi$. Hence the exact transformed
Hamiltonian has the form
\begin{equation}
 h_\eta\circ\Psi\circ\Phi_\chi
 =e+\omega\cdot I+\tfrac12I^\mathsf T\overline A I+\mathcal R_3,
 \label{eq:measure-complete-tube-normal-form}
\end{equation}
where $\mathcal R_3$ starts at total action degree three. The complete
left-hand side remains bounded by $0.05$ on the strip $1/128$ and
radius $R_1$. Cauchy's coefficient estimate in the two action
variables therefore yields, for $0<r\leq R_1/2$,
\begin{equation}
 \begin{aligned}
 \|\mathcal R_3\|_{1/128,r}
 &\leq0.05\sum_{m\geq3}(m+1)(r/R_1)^m\\
 &\leq\frac{0.5}{R_1^3}r^3=5\,10^{50}r^3.
 \end{aligned}
 \label{eq:measure-complete-cauchy-remainder}
\end{equation}
This bound includes the full magnetic cosine, all fitted quartics,
the exact chart and the exact normalization flow.

\subsection{A quantitative measure bound with all resonances included}
\label{sec:measure-ck-bound}

The action-domain dependence of KAM measure estimates is central in
\citet{biasco2018measure}. We use the fully quantified formulation
of \citet[Theorems~1 and~3, Eqs.~(4), (9), (13)--(14)]{chierchia2021measure}.
Consider a Hamiltonian $K(I)+\mathcal R(I,t)$ on
$B_R(I_c)\times\mathbb T_s^2$, with
$K=e+\omega\cdot I+\tfrac12I^\mathsf TAI$, real symmetric invertible
$A$, and $\|\mathcal R\|_{R,s}\leq E$.
Let $M\geq\|A\|_\infty$, $L\geq\|A^{-1}\|_\infty$,
$\vartheta=ML$, and $\epsilon=ME/\alpha^2$.
For $0<\rho<r\leq R$, $0<s_*<s\leq1$, and $\tau>1$, the sufficient
conditions are
\begin{equation}
 \alpha\leq c_0\rho/L,\qquad
 \epsilon\leq\epsilon_*:=
 \frac{(s-s_*)^{7(\tau+1)+10}}{c_*\vartheta^6}.
 \label{eq:measure-ck-smallness}
\end{equation}
Here $c_0,c_*$ are exactly the constants in their Appendix~A,
Remark~A.4, Eq.~(A.38), evaluated for $d=2$ and the chosen $\tau$.
Upper bounds on $M,L,E,c_*$, and a lower bound on $c_0$, are used
in the appropriate directions. The normalization is
$H=K+\varepsilon P$ with $\varepsilon\|P\|\leq E$.

Take the centre set in that theorem to be the singleton $\{I_c\}$.
The real domain is precisely $D=B_r(I_c)$, the inner domain is
$B_{r-\rho}(I_c)$, and the boundary layer in Theorem~3 is
$B_{r+\rho}(I_c)\setminus B_{r-\rho}(I_c)$, of area $16r\rho$.
Let $\mathcal K\subset D\times\mathbb T^2$ be the certified KAM set.
The Lipschitz estimate in their Eq.~(9) is at most $1/8$, so the
factor multiplying the bad action measure in Eq.~(14) is bounded
by $(25/16)(2\pi)^2$.

For completeness the remaining estimate includes every integer
resonance. For $m=|\ell|_1$, the strip
\[
 |\ell\cdot(\omega+AI)|<\alpha/m^\tau
\]
has intersection with $D$ of area at most
$4r\alpha/(m^\tau\|A\ell\|_\infty)$. Since
$\|A\ell\|_\infty\geq m/(2L)$ and there are $4m$ vectors on the
shell, summing over all shells gives
\[
 \operatorname{meas}(R_\alpha)\leq32r\alpha L\zeta(\tau).
\]
Counting both signs of each resonance gives a conservative upper
bound. Dividing Theorem~3's bound by
$\operatorname{meas}(D\times\mathbb T^2)=4r^2(2\pi)^2$ proves
\begin{equation}
 \frac{\operatorname{meas}(\mathcal K)}{4r^2(2\pi)^2}
 \geq1-\frac{25}{16}
 \left(\frac{4\rho}{r}+
       \frac{8\alpha L\zeta(\tau)}{r}\right).
 \label{eq:measure-ck-relative-lower-bound}
\end{equation}
The estimate therefore includes resonances of every integer order.

Apply this result to \eqref{eq:measure-complete-tube-normal-form}, with
\begin{equation}
 \begin{gathered}
 I_c=0,\quad r=w,\quad R=2w,\\
 s=1/128,\quad s_*=1/256,\quad \tau=2,\\
 M=3,\quad L=5100,\quad\vartheta=15300,\\
 E=4\,10^{51}w^3,\\
 \rho=w/500,\qquad\alpha=w/(1000L).
 \end{gathered}
 \label{eq:measure-tube-instantiation}
\end{equation}
Any specified positive $w\leq R_1/4$ satisfying
\begin{equation}
 4\,10^{57}M L^2w
 \leq \frac{256^{-31}}{c_*\vartheta^6}
 \label{eq:measure-explicit-w-condition}
\end{equation}
passes the smallness condition. The first condition in
\eqref{eq:measure-ck-smallness} follows already from $c_0>1$.
For every such fixed $w$, define the physical canonical tube
\begin{equation}
 U_w=(\Psi\circ\Phi_\chi)
       (\mathbb T^2\times B_w(0)).
 \label{eq:measure-physical-tube}
\end{equation}
It is an injective symplectic image of a declared domain, so its
Liouville volume is $4w^2(2\pi)^2$. Equations
\eqref{eq:measure-ck-relative-lower-bound} and
\eqref{eq:measure-tube-instantiation} give the explicit bound
\begin{equation}
 \begin{aligned}
 \frac{\mu_L(U_w\cap\mathcal K_{h_\eta})}{\mu_L(U_w)}
 &\geq1-\frac{25}{16}
       \left(\frac4{500}+\frac8{1000}\frac{\pi^2}{6}\right)\\
 &>0.9669383.
 \end{aligned}
 \label{eq:measure-fixed-tube-percentage}
\end{equation}
The numerical certificate takes the exact value $w=10^{-220}$.
It encloses $\epsilon_*$ above $7.6577\,10^{-155}$ and the
smallness ratio $\epsilon/\epsilon_*$ below $0.40759$, verifying
\eqref{eq:measure-explicit-w-condition} with outward rounding.

The transverse size of this particular tube is also bounded. For any
real point of $U_w$, write $(t',I')=\Phi_\chi(t,I)$. The same flow
estimate on the smaller action domain gives $|I'|_\infty<2w$.
Comparison with the central point $k(t')$, rather than with a fixed
phase label before the flow, eliminates the irrelevant phase shift.
For real $q>-1$, $|\sqrt{1+q}-1|
=|q|/(1+\sqrt{1+q})\leq|q|$; hence
\begin{equation}
 \begin{aligned}
 \|\Psi(t',I')-k(t')\|_\infty
 &\leq\frac{0.100007\,(1.014)\,2w}{0.005-3\,10^{-6}}\\
 &<41w=4.1\,10^{-219}.
 \end{aligned}
 \label{eq:measure-tube-transverse-width}
\end{equation}
This is a bound in the dimensionless Cartesian canonical variables;
the SI displacement bound follows by applying the fixed modal
reconstruction matrix. The millimetre value describes the amplitude
of the central motion, not the thickness in
\eqref{eq:measure-tube-transverse-width}.

The fraction concerns the conservative Hamiltonian on a four-dimensional
symplectic tube. The constant SI scaling cancels from the ratio.
As quantified in Section~\ref{sec:measure-fixed-tube}, the transverse width
is far below operational resolution and the available embedding-error
guarantee. The exact set is mathematically defined, but membership cannot
be resolved from the computed torus approximation. The lift to the damped
closed loop carries only the low-mode parameter measure specified in
Section~\ref{sec:measure-definition}.
